\section{Method}

CEMS adapts criterion-guided text scoring to the continuous rating task defined above. Here, a rubric consists of a metaphor-related dimension definition and a corresponding rating question. The LLM uses these criteria and rated examples to assign a score. Two additional predictors estimate human ratings from semantic neighbors and text embeddings. The LLM and BGE encoder remain frozen in the main configuration. Human ratings supply prompt examples and KNN neighbor labels, train the Ridge regressor, and determine the fusion weights and intercept.

Figure~\ref{fig:cems} shows how we fit the fusion weights on training data and use them to predict ratings for test items.



\subsection{Rubric-guided LLM score}

For each item and dimension, Qwen3-14B-AWQ receives the expression, the dimension definition, a rating question, and five examples retrieved from the training set and rated on that dimension. It returns a continuous score on a working scale of $[1,5]$. Retrieved examples exclude the target item and its group. Fusion uses the score; confidence and free-form rationales are not input features.

The LLM also proposes a source--target mapping and quotes evidence from the expression or available context. A validator checks whether each quote matches the stated source span. This check verifies the location of the quoted text. We test a version without the mapping prompt in an ablation.

\subsection{Semantic predictors}

We encode each expression with the frozen BAAI/bge-small-zh-v1.5 encoder \citep{xiao2024cpack}. For each dimension, BGE-KNN averages the human ratings of the five nearest training items, weighted by similarity. BGE-Ridge uses regularized linear regression to predict human means from normalized embeddings. Both use labels from the current training fold. 

\subsection{Fitting the fusion weights}

Let $\mathbf{p}_{id}=(q_{id},k_{id},r_{id})$ denote the LLM, KNN, and Ridge scores for item $i$ and dimension $d$. We learn dimension-specific nonnegative weights $\mathbf{w}_d$ that sum to one and an intercept $\beta_d$:
\begin{equation}
\hat{y}_{id}=\operatorname{clip}_{[1,5]}
\left(\mathbf{w}_d^\top\mathbf{p}_{id}+\beta_d\right).
\label{eq:fusion}
\end{equation}
We select weights using predictions generated by four-fold cross-validation within the outer training set. Each inner fold is held out in turn: prompt examples and KNN neighbors come from the remaining inner folds, and Ridge is fitted on those folds. The frozen LLM and encoder are reused throughout. Pooling the held-out predictions gives one out-of-fold (OOF) score from each branch for every outer-training item. We compare combinations of these scores against the corresponding human means using MAE. For each weight vector on the 0.05 simplex grid, we set the intercept to the median training residual, clipped to $[-0.5,0.5]$, and evaluate the resulting clipped predictions. We select the candidate with the lowest OOF MAE. We then refit the semantic predictors on the full outer training set and use the fixed fusion parameters on the test fold. Human labels in the outer test fold are never used for retrieval, regression, or calibration.
